L'obiettivo di questo esperimento è misurare la vita media dei muoni.  Per raggiungere questo scopo, sono stati utilizzati tre rivelatori messi in cascata con un attenuatore per rallentare e fermare i muoni che provengono da interazioni nell'alta atmosfera con i raggi cosmici. Un aspetto fondamentale per il successo di questa misura è il concetto di coincidenza. Infatti attraverso il circuito visibile schematicamente nella figura \ref{figura:circuito} è stato possibile selezionare i muoni che penetravano nell'attenuatore e successivamente, misurare il tempo di decadimento ricevendo un segnale di stop attraverso il passaggio del positrone nei rivelatori S1 e S2. Due misure distinte sono state eseguite in diverse configurazioni sperimentali. La prima misura è stata effettuata senza l'inserimento di attenuatori in piombo nella parte superiore della struttura che sostiene i tre rivelatori. La seconda misura, invece, è stata realizzata introducendo attenuatori in piombo, con l'obiettivo di ridurre la presenza di muoni negativi.
\begin{itemize}
    \item $\tau_1=(2.0407\pm0.1724)\mu\unit{s}$;
    \item$\tau_2=(1.8857\pm0.0989)\mu\unit{s}$.
\end{itemize}
Per la prima misura, la discrepanza rispetto al valore teorico è stata calcolata in termini di $\sigma$, risultando in uno scarto di $0.9\sigma$. Questo, insieme ai risultati del test del $\chi^2$ suggerisce una buona aderenza ai modelli teorici e una corretta esecuzione dell'esperimento

La seconda misura, tuttavia, ha mostrato una discrepanza di $3.15\sigma$ rispetto al valore teorico. Questo scarto significativo potrebbe essere indicativo di errori sistematici legati alla strumentazione, che emergono solo durante sessioni di acquisizione dati prolungate. La raccolta di dati continuativa senza interruzioni, contrariamente alla prima misura, potrebbe aver esposto l'apparato a variazioni non compensate, potenzialmente alterando la misura della vita media dei muoni.

Inoltre, è fondamentale considerare che non è stato eseguito uno studio dettagliato dei segnali di fondo che dovrebbero comunque essere ridotti dal discriminatore posto prima degli ingressi dello stadio AND di start. La mancata considerazione di questi segnali potrebbe aver contribuito a distorcere i risultati, introducendo ulteriori incertezze sistematiche non contabilizzate nella misura.